% =====================================================================
%  Chapter 6 — Discussion
%  Input with:  % =====================================================================
%  Chapter 6 — Discussion
%  Input with:  % =====================================================================
%  Chapter 6 — Discussion
%  Input with:  % =====================================================================
%  Chapter 6 — Discussion
%  Input with:  \input{chapters/06_discussion}
%  No tables or figures. Cites literature already in references.bib
%  plus four new entries (see CH6_NOTES.md).
% =====================================================================

\section{Discussion}
\label{sec:discussion}

\subsection{The size of what is left over}
\label{sec:whats-left}

Chapter~\ref{sec:premium} ended with a number and a boundary. The number is roughly 0.96 children per woman: the difference in period fertility between Central Asia and the other ten post-Soviet states that does not move together with income, urbanisation, remittances or child survival. The boundary is that this thesis cannot say what the number is made of.

It is worth holding on to the magnitude, because it is easy to read a regression coefficient as a small technical quantity. Nearly one additional child per woman, sustained across twenty-four years, in countries that shared an administrative system, a language of government, a school curriculum, a health service and a family-policy regime until 1991, is not a residual. It is larger than the entire range of variation the ten comparison countries displayed among themselves over the same period. Uzbekistan and Belarus began the period 1.4 children apart and ended it, in 2023, 2.3 apart, having been governed by the same state within living memory.

What follows sets out four channels that the literature identifies and this design cannot test, then looks at the two countries whose recent movements are largest, and closes with what the analysis cannot support.

\subsection{Four channels this design cannot test}
\label{sec:channels}

\paragraph{Marriage.} The most direct candidate is nuptiality, and the cross-section leaves a visible trace of it: female singulate mean age at marriage correlates with mean fertility at $-0.568$. The underlying figures are stark. Central Asian countries average 21.5 years, the other ten 26.3. The three Baltic states sit between 30.4 and 33.6 years, ten to thirteen years later than Tajikistan at 20.7.

The mechanism behind this is not simply that women marry younger. \textcite{dommaraju2008} show that across Kazakhstan, Kyrgyzstan and Uzbekistan marriage remained near-universal through the post-Soviet decade even as age at first marriage rose and marriage rates among younger cohorts fell considerably. Universality and timing move separately, and it is universality that matters most for aggregate fertility in societies where childbearing happens inside unions. Their evidence on education is equally relevant: in Kyrgyzstan, women with secondary schooling or less had a fifty per cent chance of still being unmarried at twenty, against ninety per cent for those with higher education. Education delays marriage sharply, and the female schooling variable in the cross-section is presumably picking up some of this.

Azerbaijan again sits with the comparison bloc rather than with Central Asia: its age at marriage is 24.2 years, close to Armenia's 24.1 and Russia's 24.9.

\paragraph{Ethnic composition.} The second channel is the one this design is furthest from being able to see, since the panel has no ethnic dimension at all. \textcite{spoorenberg2013} documents a persistent gap of about one child per woman between women of the titular ethnicities and women of European origin, in all five Central Asian republics and across multiple data sources and estimation methods. Because emigration after 1991 changed the ethnic composition of these populations substantially, that differential translates directly into national aggregates.

\textcite{spoorenberg2015} quantifies it. Holding ethnic-specific fertility fixed at its pre-increase level and letting only population composition change, the simulation reproduces 73 per cent of Kazakhstan's recorded fertility increase and 89 per cent of Kyrgyzstan's. For Tajikistan, Turkmenistan and Uzbekistan the composition effect accounts for most of the increase or stagnation, though those estimates rest on extrapolated rather than observed composition and should be treated more cautiously.

This is the single most important gap between what this thesis measures and what the literature can show. A large share of the Central Asian fertility rise, possibly the majority of it, may be a composition effect operating within countries, invisible to a design that observes only national aggregates. The Central Asia indicator absorbs it silently along with everything else.

There is a complication, though, and it cuts in a useful direction. \textcite{nedoluzhko2012} compares members of the same ethnic groups living in different Central Asian states, which separates ethnicity from national context. Ethnicity matters, but less than where people live: Tajiks have about three per cent higher fertility than Uzbeks, while living in Tajikistan rather than Uzbekistan is associated with twenty-one per cent higher fertility, holding ethnicity and minority status constant. Minority status itself has only a modest negative association. If national context dominates ethnicity even within a region as culturally continuous as Central Asia, then a country-level indicator is not as crude an instrument as it first appears, though it still cannot say which national characteristics are doing the work.

\paragraph{Religiosity, gender norms, and a generational reading.} The cross-section found mean fertility correlating with Muslim population share at $+0.880$ and then found that correlation empirically inseparable from regional identity. \textcite{kumo2024} suggest why a simple religion-to-fertility reading would have been wrong anyway.

Using World Values Survey data for five post-Soviet Muslim-majority countries against seven Muslim-majority countries outside the former Soviet Union, they model the relationship as a two-stage chain: religiosity shapes conservative gender-role beliefs, and those beliefs shape family size. Both links hold in the non-Soviet comparison group. In the post-Soviet countries the first link holds but the second does not: more conservative gender beliefs are not associated with more children. They read this as Soviet institutional legacy: decades of policy promoting female labour-force participation and formal gender equality left structures that keep conservative belief from translating into childbearing behaviour.

That finding offers one interpretation of Azerbaijan, which the cross-section could not supply. A country can be 95 per cent Muslim and have fertility of 1.88 if the institutional layer between belief and behaviour is intact.

The part of their argument most relevant to Chapter~\ref{sec:trends} is generational. When they split the post-Soviet sample by whether respondents spent their formative years under Soviet rule, the dampening holds for those who did and weakens for those who did not: among the younger group, religiosity shapes gender beliefs more strongly and those beliefs are associated with fertility. If that pattern is real, it has an implication for the timing this thesis documents. Women aged 20 to 34 in 2020 were born between 1986 and 2000, and their formative years fell almost entirely after 1991. The cohort now at peak childbearing across Central Asia is the first for which the Soviet institutional dampening was largely absent.

This is a conjecture, not a result. It is consistent with a divergence that begins around 2016--17 rather than earlier, and it is consistent with the direction of \textcite{kumo2024}'s cohort split, but nothing in this thesis tests it. It would need individual-level data with birth cohort, religiosity and gender attitudes measured together, which is precisely the sort of evidence a fourteen-country macro panel cannot substitute for.

\paragraph{Migration, remittances, and expectations.} The fourth channel is where this thesis has its one positive interaction and its most fragile result. The estimated association between remittances and fertility is more positive inside Central Asia than outside it, by 0.0157 per percentage point of GDP, against a broader post-communist literature in which the remittance--fertility relationship is generally negative \parencite{tokhirov2024}. The estimate depends on Tajikistan and Kyrgyzstan jointly. Moldova, as remittance-dependent as Kyrgyzstan and holding the sample's most negative residual, shows that remittance dependence alone produces nothing.

What might distinguish the Central Asian cases is not the transfers but what accompanies them. \textcite{agadjanian2022} find in Kyrgyzstan that completed fertility tracks group-level normative propensities while fertility \emph{desires} track something different: how a group perceives its own position and prospects. Women expecting their ethnic group's situation to improve were more likely to want another child, independent of how many they already had. Ethnic homophily in personal networks worked in the same direction. Their sharpest illustration is Uzbeks in Kyrgyzstan, a group conventional cultural reasoning would place among the most pronatalist, who instead report notably weak desires for additional children, a pattern they attribute to vulnerable societal positioning rather than to culture.

Expectations, in other words, may matter as much as income. That is a mechanism a macroeconomic panel is poorly equipped to detect, since it has no obvious annual indicator.

\subsection{Uzbekistan after 2017}
\label{sec:uzbekistan}

The largest single movement in the panel deserves separate treatment, because there is direct survey evidence about it and because it bears on the tempo question raised in Section~\ref{sec:turning-point}.

Uzbekistan's fertility rose from 2.57 in 2017 to 3.50 in 2023 in the World Population Prospects series. National figures show the same movement: annual live births increased by 27 per cent between 2017 and 2021, from about 716{,}000 to 905{,}000, with the total fertility rate rising from 2.42 to 3.17 \parencite{unfpa2023}. What makes this hard to dismiss as a compositional artefact is that it happened while the population at risk was shrinking. The number of women aged 20 to 24, a group that accounts for roughly forty per cent of Uzbek fertility, fell by 15 per cent over the same five years. Fewer women in the highest-fertility ages, and substantially more births.

The survey conducted to explain this reached conclusions worth taking seriously, partly for what it ruled out. Reduced access to contraception had been the obvious suspect; it was tested directly and not supported, and reported contraceptive use rose with the number of children a woman already had rather than falling. The factors respondents did identify were improved household finances, greater freedom in religious matters, and increased confidence in the future, with the first and last endorsed by close to ninety per cent \parencite{unfpa2023}.

That last item connects to \textcite{agadjanian2022}'s finding that expectations about the future predict fertility desires, and the two studies reach it independently, in different countries, with different methods. The survey adds a further nuance the cross-section could never have detected: respondents' own religiosity was not associated with their reproductive choices, but the religiosity of the household they lived in was. In this survey, reproductive orientations were more clearly associated with the religiosity of the household than with the respondent's own reported religiosity, which parallels, without being the same measured construct as, the group normative propensities of \textcite{agadjanian2022}.

Two cautions. The country has had no census since 1989, so all of this rests on surveys and registration rather than on an enumerated population. And the survey itself flags a warning sign: women born between 1997 and 2003 expect fewer children than they had originally planned, which the authors read as an early indication that the rise will not continue. The Uzbek surge is evidence that the Central Asian difference is live rather than residual, but it is not evidence that it will persist.

\subsection{Why the comparison bloc fell}
\label{sec:rest-fell}

Chapter~\ref{sec:trends} showed that roughly seventy per cent of the post-2016 widening came from the other ten countries declining rather than from Central Asia rising. A discussion that treated the divergence purely as a Central Asian phenomenon would therefore be explaining the smaller share of it.

Russia's trajectory illustrates the general shape. \textcite{zakharov2024} shows that the period rate began rising in 2000, six years before the maternity-capital decree that is often credited with it, and that the policy's main visible effect was to shorten intervals between births rather than to raise completed family size. Period fertility rose while births were being brought forward, then fell when that could no longer continue: by 2019 the rate was back to 1.50, essentially its level before the policy began. He notes that a comparable episode in the 1980s produced the same inflation followed by the same correction.

Something like this may have operated across the comparison bloc: postponement slowing during the recovery years, inflating period rates through the 2000s and early 2010s, then exhausting itself. On that reading the post-2016 decline is partly the unwinding of a timing effect rather than a fresh collapse in intended family size. \textcite{grigoras2019} documents postponement and recuperation patterns across Moldova, Russia, Belarus and the Baltic states consistent with this. \textcite{billingsley2010} cautions that postponement theory does not fit the region uniformly, and that when postponement begins from an age as young as the post-communist average it need not translate into forgone births at all.

The decline is also the less robust half of the divergence. Section~\ref{sec:desc-robust} shows that removing the 2022 and 2023 observations for all countries leaves Central Asian internal convergence intact but weakens the rise in dispersion among the other ten below conventional significance. This thesis observes the decline and cannot decompose it. That a substantial part of the widening difference may reflect timing on the low-fertility side rather than family size on the high-fertility side is the strongest reason for caution about the headline descriptive finding.

\subsection{Limitations}
\label{sec:limitations}

Several constraints have been noted where they arose; they are collected here.

The dependent variable throughout is the period total fertility rate, which confounds the number of births with their timing \parencite{bongaarts1998}. Tempo effects are documented on both sides of the comparison: in Central Asia by \textcite{spoorenberg2015}, who finds Kyrgyzstan's rise driven largely by reduced postponement while Kazakhstan's involved a genuine quantum increase, and in the European successor states by \textcite{grigoras2019}. No result here should be read as a statement about completed family size.

The sample is fourteen countries, four of them treated. This limits every part of the analysis: it is why the cross-section cannot separate Muslim share from regional identity, why the period interactions cannot distinguish widening from stability, and why inference required a wild-cluster bootstrap rather than conventional clustered standard errors.

The design has no ethnic, parity or cohort dimension. Given that ethnic composition may account for the majority of the Central Asian fertility increase \parencite{spoorenberg2015}, this is the most consequential omission, and it is a limitation of the data rather than of the specification.

Measurement is uneven. Eleven of the fourteen 2023 values are projections rather than retrospective estimates, and the singulate mean age at marriage is assembled from three sources with different marital-status definitions and observation years spanning 2012 to 2022, with a known bias direction that runs against the pattern of interest.

Finally, and most fundamentally, nothing here identifies a cause. The controls are plausibly consequences of fertility regimes as well as correlates of them, and the regional indicator is a label for a bundle rather than a variable. What the analysis supports is a claim about the size, stability and robustness of an association, and a claim about what that association is \emph{not} accounted for by. The step from there to explanation is one the micro-level literature can take and this design cannot.

%  No tables or figures. Cites literature already in references.bib
%  plus four new entries (see CH6_NOTES.md).
% =====================================================================

\section{Discussion}
\label{sec:discussion}

\subsection{The size of what is left over}
\label{sec:whats-left}

Chapter~\ref{sec:premium} ended with a number and a boundary. The number is roughly 0.96 children per woman: the difference in period fertility between Central Asia and the other ten post-Soviet states that does not move together with income, urbanisation, remittances or child survival. The boundary is that this thesis cannot say what the number is made of.

It is worth holding on to the magnitude, because it is easy to read a regression coefficient as a small technical quantity. Nearly one additional child per woman, sustained across twenty-four years, in countries that shared an administrative system, a language of government, a school curriculum, a health service and a family-policy regime until 1991, is not a residual. It is larger than the entire range of variation the ten comparison countries displayed among themselves over the same period. Uzbekistan and Belarus began the period 1.4 children apart and ended it, in 2023, 2.3 apart, having been governed by the same state within living memory.

What follows sets out four channels that the literature identifies and this design cannot test, then looks at the two countries whose recent movements are largest, and closes with what the analysis cannot support.

\subsection{Four channels this design cannot test}
\label{sec:channels}

\paragraph{Marriage.} The most direct candidate is nuptiality, and the cross-section leaves a visible trace of it: female singulate mean age at marriage correlates with mean fertility at $-0.568$. The underlying figures are stark. Central Asian countries average 21.5 years, the other ten 26.3. The three Baltic states sit between 30.4 and 33.6 years, ten to thirteen years later than Tajikistan at 20.7.

The mechanism behind this is not simply that women marry younger. \textcite{dommaraju2008} show that across Kazakhstan, Kyrgyzstan and Uzbekistan marriage remained near-universal through the post-Soviet decade even as age at first marriage rose and marriage rates among younger cohorts fell considerably. Universality and timing move separately, and it is universality that matters most for aggregate fertility in societies where childbearing happens inside unions. Their evidence on education is equally relevant: in Kyrgyzstan, women with secondary schooling or less had a fifty per cent chance of still being unmarried at twenty, against ninety per cent for those with higher education. Education delays marriage sharply, and the female schooling variable in the cross-section is presumably picking up some of this.

Azerbaijan again sits with the comparison bloc rather than with Central Asia: its age at marriage is 24.2 years, close to Armenia's 24.1 and Russia's 24.9.

\paragraph{Ethnic composition.} The second channel is the one this design is furthest from being able to see, since the panel has no ethnic dimension at all. \textcite{spoorenberg2013} documents a persistent gap of about one child per woman between women of the titular ethnicities and women of European origin, in all five Central Asian republics and across multiple data sources and estimation methods. Because emigration after 1991 changed the ethnic composition of these populations substantially, that differential translates directly into national aggregates.

\textcite{spoorenberg2015} quantifies it. Holding ethnic-specific fertility fixed at its pre-increase level and letting only population composition change, the simulation reproduces 73 per cent of Kazakhstan's recorded fertility increase and 89 per cent of Kyrgyzstan's. For Tajikistan, Turkmenistan and Uzbekistan the composition effect accounts for most of the increase or stagnation, though those estimates rest on extrapolated rather than observed composition and should be treated more cautiously.

This is the single most important gap between what this thesis measures and what the literature can show. A large share of the Central Asian fertility rise, possibly the majority of it, may be a composition effect operating within countries, invisible to a design that observes only national aggregates. The Central Asia indicator absorbs it silently along with everything else.

There is a complication, though, and it cuts in a useful direction. \textcite{nedoluzhko2012} compares members of the same ethnic groups living in different Central Asian states, which separates ethnicity from national context. Ethnicity matters, but less than where people live: Tajiks have about three per cent higher fertility than Uzbeks, while living in Tajikistan rather than Uzbekistan is associated with twenty-one per cent higher fertility, holding ethnicity and minority status constant. Minority status itself has only a modest negative association. If national context dominates ethnicity even within a region as culturally continuous as Central Asia, then a country-level indicator is not as crude an instrument as it first appears, though it still cannot say which national characteristics are doing the work.

\paragraph{Religiosity, gender norms, and a generational reading.} The cross-section found mean fertility correlating with Muslim population share at $+0.880$ and then found that correlation empirically inseparable from regional identity. \textcite{kumo2024} suggest why a simple religion-to-fertility reading would have been wrong anyway.

Using World Values Survey data for five post-Soviet Muslim-majority countries against seven Muslim-majority countries outside the former Soviet Union, they model the relationship as a two-stage chain: religiosity shapes conservative gender-role beliefs, and those beliefs shape family size. Both links hold in the non-Soviet comparison group. In the post-Soviet countries the first link holds but the second does not: more conservative gender beliefs are not associated with more children. They read this as Soviet institutional legacy: decades of policy promoting female labour-force participation and formal gender equality left structures that keep conservative belief from translating into childbearing behaviour.

That finding offers one interpretation of Azerbaijan, which the cross-section could not supply. A country can be 95 per cent Muslim and have fertility of 1.88 if the institutional layer between belief and behaviour is intact.

The part of their argument most relevant to Chapter~\ref{sec:trends} is generational. When they split the post-Soviet sample by whether respondents spent their formative years under Soviet rule, the dampening holds for those who did and weakens for those who did not: among the younger group, religiosity shapes gender beliefs more strongly and those beliefs are associated with fertility. If that pattern is real, it has an implication for the timing this thesis documents. Women aged 20 to 34 in 2020 were born between 1986 and 2000, and their formative years fell almost entirely after 1991. The cohort now at peak childbearing across Central Asia is the first for which the Soviet institutional dampening was largely absent.

This is a conjecture, not a result. It is consistent with a divergence that begins around 2016--17 rather than earlier, and it is consistent with the direction of \textcite{kumo2024}'s cohort split, but nothing in this thesis tests it. It would need individual-level data with birth cohort, religiosity and gender attitudes measured together, which is precisely the sort of evidence a fourteen-country macro panel cannot substitute for.

\paragraph{Migration, remittances, and expectations.} The fourth channel is where this thesis has its one positive interaction and its most fragile result. The estimated association between remittances and fertility is more positive inside Central Asia than outside it, by 0.0157 per percentage point of GDP, against a broader post-communist literature in which the remittance--fertility relationship is generally negative \parencite{tokhirov2024}. The estimate depends on Tajikistan and Kyrgyzstan jointly. Moldova, as remittance-dependent as Kyrgyzstan and holding the sample's most negative residual, shows that remittance dependence alone produces nothing.

What might distinguish the Central Asian cases is not the transfers but what accompanies them. \textcite{agadjanian2022} find in Kyrgyzstan that completed fertility tracks group-level normative propensities while fertility \emph{desires} track something different: how a group perceives its own position and prospects. Women expecting their ethnic group's situation to improve were more likely to want another child, independent of how many they already had. Ethnic homophily in personal networks worked in the same direction. Their sharpest illustration is Uzbeks in Kyrgyzstan, a group conventional cultural reasoning would place among the most pronatalist, who instead report notably weak desires for additional children, a pattern they attribute to vulnerable societal positioning rather than to culture.

Expectations, in other words, may matter as much as income. That is a mechanism a macroeconomic panel is poorly equipped to detect, since it has no obvious annual indicator.

\subsection{Uzbekistan after 2017}
\label{sec:uzbekistan}

The largest single movement in the panel deserves separate treatment, because there is direct survey evidence about it and because it bears on the tempo question raised in Section~\ref{sec:turning-point}.

Uzbekistan's fertility rose from 2.57 in 2017 to 3.50 in 2023 in the World Population Prospects series. National figures show the same movement: annual live births increased by 27 per cent between 2017 and 2021, from about 716{,}000 to 905{,}000, with the total fertility rate rising from 2.42 to 3.17 \parencite{unfpa2023}. What makes this hard to dismiss as a compositional artefact is that it happened while the population at risk was shrinking. The number of women aged 20 to 24, a group that accounts for roughly forty per cent of Uzbek fertility, fell by 15 per cent over the same five years. Fewer women in the highest-fertility ages, and substantially more births.

The survey conducted to explain this reached conclusions worth taking seriously, partly for what it ruled out. Reduced access to contraception had been the obvious suspect; it was tested directly and not supported, and reported contraceptive use rose with the number of children a woman already had rather than falling. The factors respondents did identify were improved household finances, greater freedom in religious matters, and increased confidence in the future, with the first and last endorsed by close to ninety per cent \parencite{unfpa2023}.

That last item connects to \textcite{agadjanian2022}'s finding that expectations about the future predict fertility desires, and the two studies reach it independently, in different countries, with different methods. The survey adds a further nuance the cross-section could never have detected: respondents' own religiosity was not associated with their reproductive choices, but the religiosity of the household they lived in was. In this survey, reproductive orientations were more clearly associated with the religiosity of the household than with the respondent's own reported religiosity, which parallels, without being the same measured construct as, the group normative propensities of \textcite{agadjanian2022}.

Two cautions. The country has had no census since 1989, so all of this rests on surveys and registration rather than on an enumerated population. And the survey itself flags a warning sign: women born between 1997 and 2003 expect fewer children than they had originally planned, which the authors read as an early indication that the rise will not continue. The Uzbek surge is evidence that the Central Asian difference is live rather than residual, but it is not evidence that it will persist.

\subsection{Why the comparison bloc fell}
\label{sec:rest-fell}

Chapter~\ref{sec:trends} showed that roughly seventy per cent of the post-2016 widening came from the other ten countries declining rather than from Central Asia rising. A discussion that treated the divergence purely as a Central Asian phenomenon would therefore be explaining the smaller share of it.

Russia's trajectory illustrates the general shape. \textcite{zakharov2024} shows that the period rate began rising in 2000, six years before the maternity-capital decree that is often credited with it, and that the policy's main visible effect was to shorten intervals between births rather than to raise completed family size. Period fertility rose while births were being brought forward, then fell when that could no longer continue: by 2019 the rate was back to 1.50, essentially its level before the policy began. He notes that a comparable episode in the 1980s produced the same inflation followed by the same correction.

Something like this may have operated across the comparison bloc: postponement slowing during the recovery years, inflating period rates through the 2000s and early 2010s, then exhausting itself. On that reading the post-2016 decline is partly the unwinding of a timing effect rather than a fresh collapse in intended family size. \textcite{grigoras2019} documents postponement and recuperation patterns across Moldova, Russia, Belarus and the Baltic states consistent with this. \textcite{billingsley2010} cautions that postponement theory does not fit the region uniformly, and that when postponement begins from an age as young as the post-communist average it need not translate into forgone births at all.

The decline is also the less robust half of the divergence. Section~\ref{sec:desc-robust} shows that removing the 2022 and 2023 observations for all countries leaves Central Asian internal convergence intact but weakens the rise in dispersion among the other ten below conventional significance. This thesis observes the decline and cannot decompose it. That a substantial part of the widening difference may reflect timing on the low-fertility side rather than family size on the high-fertility side is the strongest reason for caution about the headline descriptive finding.

\subsection{Limitations}
\label{sec:limitations}

Several constraints have been noted where they arose; they are collected here.

The dependent variable throughout is the period total fertility rate, which confounds the number of births with their timing \parencite{bongaarts1998}. Tempo effects are documented on both sides of the comparison: in Central Asia by \textcite{spoorenberg2015}, who finds Kyrgyzstan's rise driven largely by reduced postponement while Kazakhstan's involved a genuine quantum increase, and in the European successor states by \textcite{grigoras2019}. No result here should be read as a statement about completed family size.

The sample is fourteen countries, four of them treated. This limits every part of the analysis: it is why the cross-section cannot separate Muslim share from regional identity, why the period interactions cannot distinguish widening from stability, and why inference required a wild-cluster bootstrap rather than conventional clustered standard errors.

The design has no ethnic, parity or cohort dimension. Given that ethnic composition may account for the majority of the Central Asian fertility increase \parencite{spoorenberg2015}, this is the most consequential omission, and it is a limitation of the data rather than of the specification.

Measurement is uneven. Eleven of the fourteen 2023 values are projections rather than retrospective estimates, and the singulate mean age at marriage is assembled from three sources with different marital-status definitions and observation years spanning 2012 to 2022, with a known bias direction that runs against the pattern of interest.

Finally, and most fundamentally, nothing here identifies a cause. The controls are plausibly consequences of fertility regimes as well as correlates of them, and the regional indicator is a label for a bundle rather than a variable. What the analysis supports is a claim about the size, stability and robustness of an association, and a claim about what that association is \emph{not} accounted for by. The step from there to explanation is one the micro-level literature can take and this design cannot.

%  No tables or figures. Cites literature already in references.bib
%  plus four new entries (see CH6_NOTES.md).
% =====================================================================

\section{Discussion}
\label{sec:discussion}

\subsection{The size of what is left over}
\label{sec:whats-left}

Chapter~\ref{sec:premium} ended with a number and a boundary. The number is roughly 0.96 children per woman: the difference in period fertility between Central Asia and the other ten post-Soviet states that does not move together with income, urbanisation, remittances or child survival. The boundary is that this thesis cannot say what the number is made of.

It is worth holding on to the magnitude, because it is easy to read a regression coefficient as a small technical quantity. Nearly one additional child per woman, sustained across twenty-four years, in countries that shared an administrative system, a language of government, a school curriculum, a health service and a family-policy regime until 1991, is not a residual. It is larger than the entire range of variation the ten comparison countries displayed among themselves over the same period. Uzbekistan and Belarus began the period 1.4 children apart and ended it, in 2023, 2.3 apart, having been governed by the same state within living memory.

What follows sets out four channels that the literature identifies and this design cannot test, then looks at the two countries whose recent movements are largest, and closes with what the analysis cannot support.

\subsection{Four channels this design cannot test}
\label{sec:channels}

\paragraph{Marriage.} The most direct candidate is nuptiality, and the cross-section leaves a visible trace of it: female singulate mean age at marriage correlates with mean fertility at $-0.568$. The underlying figures are stark. Central Asian countries average 21.5 years, the other ten 26.3. The three Baltic states sit between 30.4 and 33.6 years, ten to thirteen years later than Tajikistan at 20.7.

The mechanism behind this is not simply that women marry younger. \textcite{dommaraju2008} show that across Kazakhstan, Kyrgyzstan and Uzbekistan marriage remained near-universal through the post-Soviet decade even as age at first marriage rose and marriage rates among younger cohorts fell considerably. Universality and timing move separately, and it is universality that matters most for aggregate fertility in societies where childbearing happens inside unions. Their evidence on education is equally relevant: in Kyrgyzstan, women with secondary schooling or less had a fifty per cent chance of still being unmarried at twenty, against ninety per cent for those with higher education. Education delays marriage sharply, and the female schooling variable in the cross-section is presumably picking up some of this.

Azerbaijan again sits with the comparison bloc rather than with Central Asia: its age at marriage is 24.2 years, close to Armenia's 24.1 and Russia's 24.9.

\paragraph{Ethnic composition.} The second channel is the one this design is furthest from being able to see, since the panel has no ethnic dimension at all. \textcite{spoorenberg2013} documents a persistent gap of about one child per woman between women of the titular ethnicities and women of European origin, in all five Central Asian republics and across multiple data sources and estimation methods. Because emigration after 1991 changed the ethnic composition of these populations substantially, that differential translates directly into national aggregates.

\textcite{spoorenberg2015} quantifies it. Holding ethnic-specific fertility fixed at its pre-increase level and letting only population composition change, the simulation reproduces 73 per cent of Kazakhstan's recorded fertility increase and 89 per cent of Kyrgyzstan's. For Tajikistan, Turkmenistan and Uzbekistan the composition effect accounts for most of the increase or stagnation, though those estimates rest on extrapolated rather than observed composition and should be treated more cautiously.

This is the single most important gap between what this thesis measures and what the literature can show. A large share of the Central Asian fertility rise, possibly the majority of it, may be a composition effect operating within countries, invisible to a design that observes only national aggregates. The Central Asia indicator absorbs it silently along with everything else.

There is a complication, though, and it cuts in a useful direction. \textcite{nedoluzhko2012} compares members of the same ethnic groups living in different Central Asian states, which separates ethnicity from national context. Ethnicity matters, but less than where people live: Tajiks have about three per cent higher fertility than Uzbeks, while living in Tajikistan rather than Uzbekistan is associated with twenty-one per cent higher fertility, holding ethnicity and minority status constant. Minority status itself has only a modest negative association. If national context dominates ethnicity even within a region as culturally continuous as Central Asia, then a country-level indicator is not as crude an instrument as it first appears, though it still cannot say which national characteristics are doing the work.

\paragraph{Religiosity, gender norms, and a generational reading.} The cross-section found mean fertility correlating with Muslim population share at $+0.880$ and then found that correlation empirically inseparable from regional identity. \textcite{kumo2024} suggest why a simple religion-to-fertility reading would have been wrong anyway.

Using World Values Survey data for five post-Soviet Muslim-majority countries against seven Muslim-majority countries outside the former Soviet Union, they model the relationship as a two-stage chain: religiosity shapes conservative gender-role beliefs, and those beliefs shape family size. Both links hold in the non-Soviet comparison group. In the post-Soviet countries the first link holds but the second does not: more conservative gender beliefs are not associated with more children. They read this as Soviet institutional legacy: decades of policy promoting female labour-force participation and formal gender equality left structures that keep conservative belief from translating into childbearing behaviour.

That finding offers one interpretation of Azerbaijan, which the cross-section could not supply. A country can be 95 per cent Muslim and have fertility of 1.88 if the institutional layer between belief and behaviour is intact.

The part of their argument most relevant to Chapter~\ref{sec:trends} is generational. When they split the post-Soviet sample by whether respondents spent their formative years under Soviet rule, the dampening holds for those who did and weakens for those who did not: among the younger group, religiosity shapes gender beliefs more strongly and those beliefs are associated with fertility. If that pattern is real, it has an implication for the timing this thesis documents. Women aged 20 to 34 in 2020 were born between 1986 and 2000, and their formative years fell almost entirely after 1991. The cohort now at peak childbearing across Central Asia is the first for which the Soviet institutional dampening was largely absent.

This is a conjecture, not a result. It is consistent with a divergence that begins around 2016--17 rather than earlier, and it is consistent with the direction of \textcite{kumo2024}'s cohort split, but nothing in this thesis tests it. It would need individual-level data with birth cohort, religiosity and gender attitudes measured together, which is precisely the sort of evidence a fourteen-country macro panel cannot substitute for.

\paragraph{Migration, remittances, and expectations.} The fourth channel is where this thesis has its one positive interaction and its most fragile result. The estimated association between remittances and fertility is more positive inside Central Asia than outside it, by 0.0157 per percentage point of GDP, against a broader post-communist literature in which the remittance--fertility relationship is generally negative \parencite{tokhirov2024}. The estimate depends on Tajikistan and Kyrgyzstan jointly. Moldova, as remittance-dependent as Kyrgyzstan and holding the sample's most negative residual, shows that remittance dependence alone produces nothing.

What might distinguish the Central Asian cases is not the transfers but what accompanies them. \textcite{agadjanian2022} find in Kyrgyzstan that completed fertility tracks group-level normative propensities while fertility \emph{desires} track something different: how a group perceives its own position and prospects. Women expecting their ethnic group's situation to improve were more likely to want another child, independent of how many they already had. Ethnic homophily in personal networks worked in the same direction. Their sharpest illustration is Uzbeks in Kyrgyzstan, a group conventional cultural reasoning would place among the most pronatalist, who instead report notably weak desires for additional children, a pattern they attribute to vulnerable societal positioning rather than to culture.

Expectations, in other words, may matter as much as income. That is a mechanism a macroeconomic panel is poorly equipped to detect, since it has no obvious annual indicator.

\subsection{Uzbekistan after 2017}
\label{sec:uzbekistan}

The largest single movement in the panel deserves separate treatment, because there is direct survey evidence about it and because it bears on the tempo question raised in Section~\ref{sec:turning-point}.

Uzbekistan's fertility rose from 2.57 in 2017 to 3.50 in 2023 in the World Population Prospects series. National figures show the same movement: annual live births increased by 27 per cent between 2017 and 2021, from about 716{,}000 to 905{,}000, with the total fertility rate rising from 2.42 to 3.17 \parencite{unfpa2023}. What makes this hard to dismiss as a compositional artefact is that it happened while the population at risk was shrinking. The number of women aged 20 to 24, a group that accounts for roughly forty per cent of Uzbek fertility, fell by 15 per cent over the same five years. Fewer women in the highest-fertility ages, and substantially more births.

The survey conducted to explain this reached conclusions worth taking seriously, partly for what it ruled out. Reduced access to contraception had been the obvious suspect; it was tested directly and not supported, and reported contraceptive use rose with the number of children a woman already had rather than falling. The factors respondents did identify were improved household finances, greater freedom in religious matters, and increased confidence in the future, with the first and last endorsed by close to ninety per cent \parencite{unfpa2023}.

That last item connects to \textcite{agadjanian2022}'s finding that expectations about the future predict fertility desires, and the two studies reach it independently, in different countries, with different methods. The survey adds a further nuance the cross-section could never have detected: respondents' own religiosity was not associated with their reproductive choices, but the religiosity of the household they lived in was. In this survey, reproductive orientations were more clearly associated with the religiosity of the household than with the respondent's own reported religiosity, which parallels, without being the same measured construct as, the group normative propensities of \textcite{agadjanian2022}.

Two cautions. The country has had no census since 1989, so all of this rests on surveys and registration rather than on an enumerated population. And the survey itself flags a warning sign: women born between 1997 and 2003 expect fewer children than they had originally planned, which the authors read as an early indication that the rise will not continue. The Uzbek surge is evidence that the Central Asian difference is live rather than residual, but it is not evidence that it will persist.

\subsection{Why the comparison bloc fell}
\label{sec:rest-fell}

Chapter~\ref{sec:trends} showed that roughly seventy per cent of the post-2016 widening came from the other ten countries declining rather than from Central Asia rising. A discussion that treated the divergence purely as a Central Asian phenomenon would therefore be explaining the smaller share of it.

Russia's trajectory illustrates the general shape. \textcite{zakharov2024} shows that the period rate began rising in 2000, six years before the maternity-capital decree that is often credited with it, and that the policy's main visible effect was to shorten intervals between births rather than to raise completed family size. Period fertility rose while births were being brought forward, then fell when that could no longer continue: by 2019 the rate was back to 1.50, essentially its level before the policy began. He notes that a comparable episode in the 1980s produced the same inflation followed by the same correction.

Something like this may have operated across the comparison bloc: postponement slowing during the recovery years, inflating period rates through the 2000s and early 2010s, then exhausting itself. On that reading the post-2016 decline is partly the unwinding of a timing effect rather than a fresh collapse in intended family size. \textcite{grigoras2019} documents postponement and recuperation patterns across Moldova, Russia, Belarus and the Baltic states consistent with this. \textcite{billingsley2010} cautions that postponement theory does not fit the region uniformly, and that when postponement begins from an age as young as the post-communist average it need not translate into forgone births at all.

The decline is also the less robust half of the divergence. Section~\ref{sec:desc-robust} shows that removing the 2022 and 2023 observations for all countries leaves Central Asian internal convergence intact but weakens the rise in dispersion among the other ten below conventional significance. This thesis observes the decline and cannot decompose it. That a substantial part of the widening difference may reflect timing on the low-fertility side rather than family size on the high-fertility side is the strongest reason for caution about the headline descriptive finding.

\subsection{Limitations}
\label{sec:limitations}

Several constraints have been noted where they arose; they are collected here.

The dependent variable throughout is the period total fertility rate, which confounds the number of births with their timing \parencite{bongaarts1998}. Tempo effects are documented on both sides of the comparison: in Central Asia by \textcite{spoorenberg2015}, who finds Kyrgyzstan's rise driven largely by reduced postponement while Kazakhstan's involved a genuine quantum increase, and in the European successor states by \textcite{grigoras2019}. No result here should be read as a statement about completed family size.

The sample is fourteen countries, four of them treated. This limits every part of the analysis: it is why the cross-section cannot separate Muslim share from regional identity, why the period interactions cannot distinguish widening from stability, and why inference required a wild-cluster bootstrap rather than conventional clustered standard errors.

The design has no ethnic, parity or cohort dimension. Given that ethnic composition may account for the majority of the Central Asian fertility increase \parencite{spoorenberg2015}, this is the most consequential omission, and it is a limitation of the data rather than of the specification.

Measurement is uneven. Eleven of the fourteen 2023 values are projections rather than retrospective estimates, and the singulate mean age at marriage is assembled from three sources with different marital-status definitions and observation years spanning 2012 to 2022, with a known bias direction that runs against the pattern of interest.

Finally, and most fundamentally, nothing here identifies a cause. The controls are plausibly consequences of fertility regimes as well as correlates of them, and the regional indicator is a label for a bundle rather than a variable. What the analysis supports is a claim about the size, stability and robustness of an association, and a claim about what that association is \emph{not} accounted for by. The step from there to explanation is one the micro-level literature can take and this design cannot.

%  No tables or figures. Cites literature already in references.bib
%  plus four new entries (see CH6_NOTES.md).
% =====================================================================

\section{Discussion}
\label{sec:discussion}

\subsection{The size of what is left over}
\label{sec:whats-left}

Chapter~\ref{sec:premium} ended with a number and a boundary. The number is roughly 0.96 children per woman: the difference in period fertility between Central Asia and the other ten post-Soviet states that does not move together with income, urbanisation, remittances or child survival. The boundary is that this thesis cannot say what the number is made of.

It is worth holding on to the magnitude, because it is easy to read a regression coefficient as a small technical quantity. Nearly one additional child per woman, sustained across twenty-four years, in countries that shared an administrative system, a language of government, a school curriculum, a health service and a family-policy regime until 1991, is not a residual. It is larger than the entire range of variation the ten comparison countries displayed among themselves over the same period. Uzbekistan and Belarus began the period 1.4 children apart and ended it, in 2023, 2.3 apart, having been governed by the same state within living memory.

What follows sets out four channels that the literature identifies and this design cannot test, then looks at the two countries whose recent movements are largest, and closes with what the analysis cannot support.

\subsection{Four channels this design cannot test}
\label{sec:channels}

\paragraph{Marriage.} The most direct candidate is nuptiality, and the cross-section leaves a visible trace of it: female singulate mean age at marriage correlates with mean fertility at $-0.568$. The underlying figures are stark. Central Asian countries average 21.5 years, the other ten 26.3. The three Baltic states sit between 30.4 and 33.6 years, ten to thirteen years later than Tajikistan at 20.7.

The mechanism behind this is not simply that women marry younger. \textcite{dommaraju2008} show that across Kazakhstan, Kyrgyzstan and Uzbekistan marriage remained near-universal through the post-Soviet decade even as age at first marriage rose and marriage rates among younger cohorts fell considerably. Universality and timing move separately, and it is universality that matters most for aggregate fertility in societies where childbearing happens inside unions. Their evidence on education is equally relevant: in Kyrgyzstan, women with secondary schooling or less had a fifty per cent chance of still being unmarried at twenty, against ninety per cent for those with higher education. Education delays marriage sharply, and the female schooling variable in the cross-section is presumably picking up some of this.

Azerbaijan again sits with the comparison bloc rather than with Central Asia: its age at marriage is 24.2 years, close to Armenia's 24.1 and Russia's 24.9.

\paragraph{Ethnic composition.} The second channel is the one this design is furthest from being able to see, since the panel has no ethnic dimension at all. \textcite{spoorenberg2013} documents a persistent gap of about one child per woman between women of the titular ethnicities and women of European origin, in all five Central Asian republics and across multiple data sources and estimation methods. Because emigration after 1991 changed the ethnic composition of these populations substantially, that differential translates directly into national aggregates.

\textcite{spoorenberg2015} quantifies it. Holding ethnic-specific fertility fixed at its pre-increase level and letting only population composition change, the simulation reproduces 73 per cent of Kazakhstan's recorded fertility increase and 89 per cent of Kyrgyzstan's. For Tajikistan, Turkmenistan and Uzbekistan the composition effect accounts for most of the increase or stagnation, though those estimates rest on extrapolated rather than observed composition and should be treated more cautiously.

This is the single most important gap between what this thesis measures and what the literature can show. A large share of the Central Asian fertility rise, possibly the majority of it, may be a composition effect operating within countries, invisible to a design that observes only national aggregates. The Central Asia indicator absorbs it silently along with everything else.

There is a complication, though, and it cuts in a useful direction. \textcite{nedoluzhko2012} compares members of the same ethnic groups living in different Central Asian states, which separates ethnicity from national context. Ethnicity matters, but less than where people live: Tajiks have about three per cent higher fertility than Uzbeks, while living in Tajikistan rather than Uzbekistan is associated with twenty-one per cent higher fertility, holding ethnicity and minority status constant. Minority status itself has only a modest negative association. If national context dominates ethnicity even within a region as culturally continuous as Central Asia, then a country-level indicator is not as crude an instrument as it first appears, though it still cannot say which national characteristics are doing the work.

\paragraph{Religiosity, gender norms, and a generational reading.} The cross-section found mean fertility correlating with Muslim population share at $+0.880$ and then found that correlation empirically inseparable from regional identity. \textcite{kumo2024} suggest why a simple religion-to-fertility reading would have been wrong anyway.

Using World Values Survey data for five post-Soviet Muslim-majority countries against seven Muslim-majority countries outside the former Soviet Union, they model the relationship as a two-stage chain: religiosity shapes conservative gender-role beliefs, and those beliefs shape family size. Both links hold in the non-Soviet comparison group. In the post-Soviet countries the first link holds but the second does not: more conservative gender beliefs are not associated with more children. They read this as Soviet institutional legacy: decades of policy promoting female labour-force participation and formal gender equality left structures that keep conservative belief from translating into childbearing behaviour.

That finding offers one interpretation of Azerbaijan, which the cross-section could not supply. A country can be 95 per cent Muslim and have fertility of 1.88 if the institutional layer between belief and behaviour is intact.

The part of their argument most relevant to Chapter~\ref{sec:trends} is generational. When they split the post-Soviet sample by whether respondents spent their formative years under Soviet rule, the dampening holds for those who did and weakens for those who did not: among the younger group, religiosity shapes gender beliefs more strongly and those beliefs are associated with fertility. If that pattern is real, it has an implication for the timing this thesis documents. Women aged 20 to 34 in 2020 were born between 1986 and 2000, and their formative years fell almost entirely after 1991. The cohort now at peak childbearing across Central Asia is the first for which the Soviet institutional dampening was largely absent.

This is a conjecture, not a result. It is consistent with a divergence that begins around 2016--17 rather than earlier, and it is consistent with the direction of \textcite{kumo2024}'s cohort split, but nothing in this thesis tests it. It would need individual-level data with birth cohort, religiosity and gender attitudes measured together, which is precisely the sort of evidence a fourteen-country macro panel cannot substitute for.

\paragraph{Migration, remittances, and expectations.} The fourth channel is where this thesis has its one positive interaction and its most fragile result. The estimated association between remittances and fertility is more positive inside Central Asia than outside it, by 0.0157 per percentage point of GDP, against a broader post-communist literature in which the remittance--fertility relationship is generally negative \parencite{tokhirov2024}. The estimate depends on Tajikistan and Kyrgyzstan jointly. Moldova, as remittance-dependent as Kyrgyzstan and holding the sample's most negative residual, shows that remittance dependence alone produces nothing.

What might distinguish the Central Asian cases is not the transfers but what accompanies them. \textcite{agadjanian2022} find in Kyrgyzstan that completed fertility tracks group-level normative propensities while fertility \emph{desires} track something different: how a group perceives its own position and prospects. Women expecting their ethnic group's situation to improve were more likely to want another child, independent of how many they already had. Ethnic homophily in personal networks worked in the same direction. Their sharpest illustration is Uzbeks in Kyrgyzstan, a group conventional cultural reasoning would place among the most pronatalist, who instead report notably weak desires for additional children, a pattern they attribute to vulnerable societal positioning rather than to culture.

Expectations, in other words, may matter as much as income. That is a mechanism a macroeconomic panel is poorly equipped to detect, since it has no obvious annual indicator.

\subsection{Uzbekistan after 2017}
\label{sec:uzbekistan}

The largest single movement in the panel deserves separate treatment, because there is direct survey evidence about it and because it bears on the tempo question raised in Section~\ref{sec:turning-point}.

Uzbekistan's fertility rose from 2.57 in 2017 to 3.50 in 2023 in the World Population Prospects series. National figures show the same movement: annual live births increased by 27 per cent between 2017 and 2021, from about 716{,}000 to 905{,}000, with the total fertility rate rising from 2.42 to 3.17 \parencite{unfpa2023}. What makes this hard to dismiss as a compositional artefact is that it happened while the population at risk was shrinking. The number of women aged 20 to 24, a group that accounts for roughly forty per cent of Uzbek fertility, fell by 15 per cent over the same five years. Fewer women in the highest-fertility ages, and substantially more births.

The survey conducted to explain this reached conclusions worth taking seriously, partly for what it ruled out. Reduced access to contraception had been the obvious suspect; it was tested directly and not supported, and reported contraceptive use rose with the number of children a woman already had rather than falling. The factors respondents did identify were improved household finances, greater freedom in religious matters, and increased confidence in the future, with the first and last endorsed by close to ninety per cent \parencite{unfpa2023}.

That last item connects to \textcite{agadjanian2022}'s finding that expectations about the future predict fertility desires, and the two studies reach it independently, in different countries, with different methods. The survey adds a further nuance the cross-section could never have detected: respondents' own religiosity was not associated with their reproductive choices, but the religiosity of the household they lived in was. In this survey, reproductive orientations were more clearly associated with the religiosity of the household than with the respondent's own reported religiosity, which parallels, without being the same measured construct as, the group normative propensities of \textcite{agadjanian2022}.

Two cautions. The country has had no census since 1989, so all of this rests on surveys and registration rather than on an enumerated population. And the survey itself flags a warning sign: women born between 1997 and 2003 expect fewer children than they had originally planned, which the authors read as an early indication that the rise will not continue. The Uzbek surge is evidence that the Central Asian difference is live rather than residual, but it is not evidence that it will persist.

\subsection{Why the comparison bloc fell}
\label{sec:rest-fell}

Chapter~\ref{sec:trends} showed that roughly seventy per cent of the post-2016 widening came from the other ten countries declining rather than from Central Asia rising. A discussion that treated the divergence purely as a Central Asian phenomenon would therefore be explaining the smaller share of it.

Russia's trajectory illustrates the general shape. \textcite{zakharov2024} shows that the period rate began rising in 2000, six years before the maternity-capital decree that is often credited with it, and that the policy's main visible effect was to shorten intervals between births rather than to raise completed family size. Period fertility rose while births were being brought forward, then fell when that could no longer continue: by 2019 the rate was back to 1.50, essentially its level before the policy began. He notes that a comparable episode in the 1980s produced the same inflation followed by the same correction.

Something like this may have operated across the comparison bloc: postponement slowing during the recovery years, inflating period rates through the 2000s and early 2010s, then exhausting itself. On that reading the post-2016 decline is partly the unwinding of a timing effect rather than a fresh collapse in intended family size. \textcite{grigoras2019} documents postponement and recuperation patterns across Moldova, Russia, Belarus and the Baltic states consistent with this. \textcite{billingsley2010} cautions that postponement theory does not fit the region uniformly, and that when postponement begins from an age as young as the post-communist average it need not translate into forgone births at all.

The decline is also the less robust half of the divergence. Section~\ref{sec:desc-robust} shows that removing the 2022 and 2023 observations for all countries leaves Central Asian internal convergence intact but weakens the rise in dispersion among the other ten below conventional significance. This thesis observes the decline and cannot decompose it. That a substantial part of the widening difference may reflect timing on the low-fertility side rather than family size on the high-fertility side is the strongest reason for caution about the headline descriptive finding.

\subsection{Limitations}
\label{sec:limitations}

Several constraints have been noted where they arose; they are collected here.

The dependent variable throughout is the period total fertility rate, which confounds the number of births with their timing \parencite{bongaarts1998}. Tempo effects are documented on both sides of the comparison: in Central Asia by \textcite{spoorenberg2015}, who finds Kyrgyzstan's rise driven largely by reduced postponement while Kazakhstan's involved a genuine quantum increase, and in the European successor states by \textcite{grigoras2019}. No result here should be read as a statement about completed family size.

The sample is fourteen countries, four of them treated. This limits every part of the analysis: it is why the cross-section cannot separate Muslim share from regional identity, why the period interactions cannot distinguish widening from stability, and why inference required a wild-cluster bootstrap rather than conventional clustered standard errors.

The design has no ethnic, parity or cohort dimension. Given that ethnic composition may account for the majority of the Central Asian fertility increase \parencite{spoorenberg2015}, this is the most consequential omission, and it is a limitation of the data rather than of the specification.

Measurement is uneven. Eleven of the fourteen 2023 values are projections rather than retrospective estimates, and the singulate mean age at marriage is assembled from three sources with different marital-status definitions and observation years spanning 2012 to 2022, with a known bias direction that runs against the pattern of interest.

Finally, and most fundamentally, nothing here identifies a cause. The controls are plausibly consequences of fertility regimes as well as correlates of them, and the regional indicator is a label for a bundle rather than a variable. What the analysis supports is a claim about the size, stability and robustness of an association, and a claim about what that association is \emph{not} accounted for by. The step from there to explanation is one the micro-level literature can take and this design cannot.
